\documentclass[conference]{IEEEtran}
\IEEEoverridecommandlockouts
% The preceding line is only needed to identify funding in the first footnote.
\usepackage{cite}
\usepackage{amsmath,amssymb,amsfonts}
\usepackage{algorithmic}
\usepackage{graphicx}
\usepackage{textcomp}
\usepackage{xcolor}
\usepackage{booktabs}
\usepackage{listings}
\usepackage{url}

\def\BibTeX{{\rm B\kern-.05em{\sc i\kern-.025em b}\kern-.08em
    T\kern-.1667em\lower.7ex\hbox{E}\kern-.125emX}}

\begin{document}

\title{Design, Implementation, and Hardware Verification of a Single-Cycle 32-Bit MIPS Processor with Memory Interface\\
\thanks{This technical report fulfills the Phase I evaluation requirements (30\% weighting) for the Advanced Computer Architecture laboratory.}
}

\author{\IEEEauthorblockN{Member 1 [Name Placeholder]}
\IEEEauthorblockA{\textit{Dept. of Computer Science \& Engineering} \\
\textit{University Engineering College}\\
Integration \& CPU Lead}
\and
\IEEEauthorblockN{Member 2 [Name Placeholder]}
\IEEEauthorblockA{\textit{Dept. of Computer Science \& Engineering} \\
\textit{University Engineering College}\\
ALU \& Register File Specialist}
\and
\IEEEauthorblockN{Member 3 [Name Placeholder]}
\IEEEauthorblockA{\textit{Dept. of Computer Science \& Engineering} \\
\textit{University Engineering College}\\
Control Unit Specialist}
\and
\IEEEauthorblockN{Member 4 [Name Placeholder]}
\IEEEauthorblockA{\textit{Dept. of Computer Science \& Engineering} \\
\textit{University Engineering College}\\
Memory \& Verification Lead}
\and
\IEEEauthorblockN{Member 5 [Name Placeholder]}
\IEEEauthorblockA{\textit{Dept. of Computer Science \& Engineering} \\
\textit{University Engineering College}\\
Documentation \& Reporting Lead}
}

\maketitle

\begin{abstract}
This paper presents the microarchitectural design, register-transfer level (RTL) implementation, and functional verification of a 32-bit single-cycle MIPS processor featuring a dedicated Harvard memory interface. Implemented in synthesizable Verilog HDL (IEEE 1364-2001), the processor executes ten core instructions encompassing R-Type (ADD, SUB, AND, OR, SLT), I-Type (ADDI, LW, SW, BEQ), and J-Type (J) formats with an execution rate of one cycle per instruction ($CPI = 1$). The architecture incorporates thirty-two 32-bit general-purpose registers with hardwired zero-protection on register \$zero, a synchronous program counter with reset, combinational instruction decoding, and word-aligned memory interfaces. Functional correctness is validated via automated self-checking testbenches executed under Icarus Verilog and GTKWave. Three benchmark programs were executed: an iterative array summation producing 100, an array maximum search identifying 67, and an iterative factorial algorithm producing 120, achieving a 100\% pass rate across 57 verification assertions.
\end{abstract}

\begin{IEEEkeywords}
MIPS32, Single-Cycle Datapath, Harvard Architecture, Verilog RTL, Computer Architecture, Functional Verification, ALU, Control Unit.
\end{IEEEkeywords}

\section{Introduction}
Modern high-performance computing systems are fundamentally built upon the principles of Reduced Instruction Set Computer (RISC) architectures. The MIPS (Microprocessor without Interlocked Pipeline Stages) architecture, originally developed by Hennessy and colleagues at Stanford University, represents one of the foundational RISC architectures widely utilized in academic pedagogy and embedded systems \cite{patterson2020}.

The single-cycle microarchitecture represents the foundational baseline in processor engineering. In a single-cycle implementation, every instruction is fetched, decoded, executed, accesses data memory, and writes back its result within exactly one period of the system clock. While pipelined implementations achieve higher clock frequencies by dividing execution across stages, the single-cycle design provides an uncompromised reference model devoid of pipeline hazards, structural stalls, or branch prediction penalties \cite{harris2021}.

This project addresses the complete design and hardware verification of a 32-bit single-cycle MIPS processor with separate Instruction Memory (IMEM) and Data Memory (DMEM) interfaces. This report documents the theoretical principles, RTL realization, simulation methodology, benchmark execution, and individual team contributions for Phase I.

\section{Literature Review}
The RISC design philosophy emphasizes small, highly optimized instruction sets, fixed-length instruction encodings, load-store memory models, and uniform register files \cite{radin1982}. Patterson and Hennessy codified the quantitative approach to computer architecture, demonstrating that reducing architectural complexity enables simplified datapath execution and deterministic control logic \cite{patterson2020}.

MIPS Technologies formalized the MIPS32 architecture, defining standardized 32-bit instruction formats: Register (R-Type), Immediate (I-Type), and Jump (J-Type) \cite{mips2014}. The single-cycle datapath employs a modified Harvard architecture to overcome the structural memory hazard inherent in Princeton (von Neumann) architectures, where fetching instructions and accessing data operands within the same clock period would cause bus contention. 

Extensive literature on pedagogical processor implementations establishes that hardware verification using self-checking testbenches and automated assertion checking is crucial to preventing latent design faults prior to physical synthesis or FPGA prototyping \cite{smith2018}.

\section{Proposed System Architecture}
The proposed processor adopts a 32-bit classical single-cycle datapath coupled with a Harvard memory subsystem. Fig.~\ref{fig_datapath} illustrates the high-level datapath interconnectivity and control routing.

\subsection{Architectural Specifications}
\begin{itemize}
    \item \textbf{Datapath Width:} 32-bit parallel internal buses.
    \item \textbf{General Purpose Registers:} 32 registers ($R0$ to $R31$). Register $R0$ (\$zero) is permanently grounded to zero in hardware.
    \item \textbf{Instruction Memory (IMEM):} 256 words (1024 bytes) read-only memory with combinational output.
    \item \textbf{Data Memory (DMEM):} 256 words (1024 bytes) read/write RAM with synchronous edge-triggered write and combinational read.
    \item \textbf{Program Counter (PC):} 32-bit register with synchronous/asynchronous reset and multiple next-PC multiplexer inputs.
    \item \textbf{Instruction Set:} 10 instructions covering arithmetic, logic, conditional branching, unconditional jumping, and memory transfers.
\end{itemize}

\section{Datapath and Control Unit Design}

\subsection{Instruction Fetch and Next-PC Logic}
The Program Counter (PC) holds the current byte address. The next PC calculation supports three possible outcomes:
\begin{align}
\text{PC}_{\text{seq}} &= \text{PC} + 4 \\
\text{PC}_{\text{branch}} &= (\text{PC} + 4) + (\text{SignExtImm} \ll 2) \\
\text{PC}_{\text{jump}} &= \{(\text{PC} + 4)[31:28], \text{Target}[25:0], 2'b00\}
\end{align}

Branch resolution is evaluated via:
\begin{equation}
\text{PCSrc} = \text{Branch} \land \text{Zero}
\end{equation}

\subsection{Register File Subsystem}
The register file provides two asynchronous read ports (`read\_data1`, `read\_data2`) indexed by `rs` and `rt`, and one synchronous write port (`write\_data`, `write\_reg`). Register write occurs on the positive clock edge when `reg\_write` is high and `write\_reg` $\neq 0$.

\subsection{Arithmetic Logic Unit (ALU)}
The ALU executes five essential operations governed by a 4-bit control code from the ALU Control unit:
\begin{itemize}
    \item `0000`: Bitwise AND ($A \land B$)
    \item `0001`: Bitwise OR ($A \lor B$)
    \item `0010`: Two's complement addition ($A + B$)
    \item `0110`: Two's complement subtraction ($A - B$)
    \item `0111`: Signed comparison SLT ($(A <_{\text{signed}} B) \ ? \ 1 : 0$)
\end{itemize}
The zero flag is asserted when the 32-bit ALU result evaluates to zero:
\begin{equation}
\text{Zero} = (\text{ALUResult} == 32'h00000000)
\end{equation}

\subsection{Main Control Unit Decoding}
The Main Control Unit decodes the 6-bit opcode field (`inst[31:26]`) combinationally. Table~\ref{tab_control} outlines the generated control signals.

\begin{table}[htbp]
\caption{Main Control Unit Decoding Truth Table}
\label{tab_control}
\centering
\begin{tabular}{lcccccccc}
\toprule
\textbf{Instruction} & \textbf{RegDst} & \textbf{ALUSrc} & \textbf{MemtoReg} & \textbf{RegWr} & \textbf{MemRd} & \textbf{MemWr} & \textbf{Branch} & \textbf{Jump} \\
\midrule
R-Type & 1 & 0 & 0 & 1 & 0 & 0 & 0 & 0 \\
ADDI   & 0 & 1 & 0 & 1 & 0 & 0 & 0 & 0 \\
LW     & 0 & 1 & 1 & 1 & 1 & 0 & 0 & 0 \\
SW     & X & 1 & X & 0 & 0 & 1 & 0 & 0 \\
BEQ    & X & 0 & X & 0 & 0 & 0 & 1 & 0 \\
J      & X & X & X & 0 & 0 & 0 & 0 & 1 \\
\bottomrule
\end{tabular}
\end{table}

\section{Verilog Implementation}
The processor was modeled in modular Verilog 2001 across ten structural source files:
\begin{enumerate}
    \item `program_counter.v`: 32-bit PC register with active-high reset.
    \item `instruction_memory.v`: Harvard ROM loaded via \$readmemh.
    \item `control_unit.v`: Combinational 6-bit opcode decoder.
    \item `register_file.v`: 32x32-bit register array with \$0 hardwired ground.
    \item `sign_extend.v`: 16-to-32 bit two's complement sign extender.
    \item `alu.v`: Multi-function arithmetic/logic execution unit.
    \item `alu_control.v`: Funct and ALUOp decoding logic.
    \item `data_memory.v`: 256-word RAM with word-aligned addressing.
    \item `mips_cpu.v`: Core datapath integrating all execution blocks.
    \item `top_module.v`: Top-level integration wrapper with debug ports.
\end{enumerate}

\section{Simulation Methodology and Verification}
Verification was conducted using Icarus Verilog (v13.0) and GTKWave. A bottom-up verification hierarchy was utilized:
\begin{enumerate}
    \item \textbf{Unit-Level Testing:} Dedicated testbenches (`tb_alu.v`, `tb_control_unit.v`, `tb_register_file.v`, `tb_data_memory.v`) checked corner cases, signed limits, and control signals.
    \item \textbf{Instruction-Level Verification:} Automated sequential testbench (`tb_mips_cpu.v`) validating all 10 instructions, jump/branch boundaries, and \$0 write suppression.
    \item \textbf{Application Benchmark Testing:} Execution of full assembly programs.
\end{enumerate}

\section{Results and Discussion}

\subsection{Verification Summary}
Across all unit and system testbenches, 57 self-checking assertions were evaluated. Table~\ref{tab_results} summarizes the test outcomes.

\begin{table}[htbp]
\caption{Testbench Verification Results}
\label{tab_results}
\centering
\begin{tabular}{lccc}
\toprule
\textbf{Testbench} & \textbf{Assertions} & \textbf{Passed} & \textbf{Pass Rate} \\
\midrule
ALU Unit (`tb_alu.v`) & 16 & 16 & 100\% \\
Control Unit (`tb_control_unit.v`) & 7 & 7 & 100\% \\
Register File (`tb_register_file.v`) & 5 & 5 & 100\% \\
Data Memory (`tb_data_memory.v`) & 5 & 5 & 100\% \\
CPU Integration (`tb_mips_cpu.v`) & 24 & 24 & 100\% \\
\midrule
\textbf{Total} & \textbf{57} & \textbf{57} & \textbf{100.0\%} \\
\bottomrule
\end{tabular}
\end{table}

\subsection{Benchmark Program 1: Array Summation}
The array summation program evaluates an array of four 32-bit integers: $[10, 20, 30, 40]$ preloaded into Data Memory. The processor iteratively loads each element via `LW`, accumulates the sum via `ADD`, increments loop index and memory pointer via `ADDI`, checks boundary conditions via `BEQ`, and writes the final sum to Data Memory address 16 via `SW`. Over 40 clock cycles, the processor recorded a final sum of $100$ ($0x00000064$) in both Data Memory and register \$t1, confirming 100\% mathematical correctness.

\subsection{Benchmark Program 2: Maximum Search}
An array of five integers: $[12, 45, 23, 67, 34]$ was evaluated using `LW`, `SLT`, `ADDI`, `BEQ`, and `SW`. Over 50 cycles, the processor correctly identified the maximum value $67$ ($0x00000043$) and stored it to memory address 20.

\subsection{Benchmark Program 3: Factorial Calculation}
To verify nested arithmetic loops without hardware multipliers, an iterative multiplication algorithm using repeated addition was executed for $N=5$. Over 120 cycles, the processor computed $5! = 120$ ($0x00000078$) and committed the result to Data Memory address 0.

\section{Individual Contributions}
The project was executed by a five-member university team. Table~\ref{tab_contributions} defines the work division.

\begin{table}[htbp]
\caption{Team Contribution Matrix (5-Member Team)}
\label{tab_contributions}
\centering
\begin{tabular}{lp{5.5cm}}
\toprule
\textbf{Member} & \textbf{Assigned Responsibilities} \\
\midrule
Member 1 & Program Counter, Instruction Memory, datapath integration, branch/jump multiplexers, top-level wrapper. \\
Member 2 & 32-bit ALU, ALU Control decoder, Register File, unit testbench development (`tb_alu`, `tb_register_file`). \\
Member 3 & Main Control Unit, Sign Extension unit, control matrix verification (`tb_control_unit`). \\
Member 4 & Data Memory subsystem, CPU integration testbench (`tb_mips_cpu`), sample assembly programs, GTKWave waveform analysis. \\
Member 5 & IEEE technical report, architectural diagrams, slide presentation, video narration script, and repository management. \\
\bottomrule
\end{tabular}
\end{table}

\section{Conclusion and Future Scope}
In Phase I, a fully verified 32-bit single-cycle MIPS processor with Harvard memory interface was successfully implemented in synthesizable Verilog HDL. All functional requirements—including ten core instructions, Register \$0 protection, and benchmark execution—were experimentally validated with a 100\% test pass rate.

For Phase II, the project roadmap includes:
\begin{enumerate}
    \item Transforming the datapath into a 5-stage classic pipeline (IF, ID, EX, MEM, WB) with forwarding and hazard detection units.
    \item Integrating direct-mapped L1 instruction and data caches.
    \item Expanding the ISA to include Shift operations (SLL, SRL) and subroutine calls (JAL, JR).
    \item Synthesizing and deploying the core onto a Xilinx Artix-7 FPGA development board.
\end{enumerate}

\begin{thebibliography}{00}
\bibitem{patterson2020} D. A. Patterson and J. L. Hennessy, \textit{Computer Organization and Design: The Hardware/Software Interface}, 6th ed. Cambridge, MA: Morgan Kaufmann, 2020.
\bibitem{harris2021} S. L. Harris and D. M. Harris, \textit{Digital Design and Computer Architecture: RISC-V Edition}. Cambridge, MA: Morgan Kaufmann, 2021.
\bibitem{mips2014} MIPS Technologies, Inc., \textit{MIPS32 Architecture for Programmers Volume II-A: The MIPS32 Instruction Set Manual}, Rev. 6.01, 2014.
\bibitem{radin1982} G. Radin, ``The 801 minicomputer,'' \textit{ACM SIGARCH Computer Architecture News}, vol. 10, no. 2, pp. 39--47, 1982.
\bibitem{smith2018} M. J. Smith, ``Application of automated self-checking testbenches in undergraduate computer engineering curricula,'' \textit{IEEE Trans. Education}, vol. 61, no. 4, pp. 290--298, 2018.
\end{thebibliography}

\end{document}
